u% -----------------------------------------------------------
\section{Obedience}
% -----------------------------------------------------------
	\label{OBEDIENCE}
	
% % ----------------------
% \subsection{See also}
% % ----------------------

% % ----------------------
% \subsection{1828 American Dictionary of the English Language} % *1828
% % ----------------------
	% \label{OBEDIENCE-1828}

	% \begin{compactitem}
		% \item
	% \end{compactitem}

% % ----------------------
% \subsection{Oxford English Dictionary} % *OED
% % ----------------------
	% \label{OBEDIENCE-OED}

	% \begin{compactitem}
		% \item[]
	% \end{compactitem}
	
% % ----------------------
% \subsection{Restoration Lexicon} % *RL *RESTORATION LEXICON
% % ----------------------
	% \label{OBEDIENCE-OED}

	% \begin{compactitem}
		% \item[]
	% \end{compactitem}

% ----------------------
\subsection{Scriptural Usage} % *SCRIPTURES
% ----------------------
	\label{OBEDIENCE-SCRIPTURES}

	% -----
	\subsubsection{Old Testament} % *OT
	% -----
		\label{OBEDIENCE-OT}
	
		% --
		\paragraph{KJV}
		% --
			\label{OBEDIENCE-OT-KJV}
	
			% \begin{tabular}{ll}
				% Total occurrences in the OT & 
			% \end{tabular}
			
			\begin{compactdesc}
				\item[{\hyperref[1SAMUEL-15-22]{1 Samuel 15:22}}] And Samuel said, Hath the LORD as great delight in burnt offerings and sacrifices, as in obeying the voice of the LORD?  Behold, to obey is better than sacrifice, and to hearken than the fat of rams.
			\end{compactdesc}
			
		% --
		\paragraph{Hebrew Bible}
		% --
			\label{OBEDIENCE-HEBREW}
		
			%
			\subparagraph{\texorpdfstring{\he{+sAma`}}{}} % שָׁמַע
			%
				\label{SHAMA} % whatever is in step bible without periods
			
				\begin{compactdesc}
					\item[HALOT 1570a, PDF 1626] \textbf{} % Use the 1994 PDF
						\begin{compactitem}
							\item[\textbf{QAL}]
							\item[1] to \textbf{hear} with the ears
							\item[1.a] seldom used absolutely
							\item[1.b] with accusative of the thing
							\item[1.c] corresponding to an Akkadian form, a witness to what he had heard, witness to a statement
							\item[1.d] to hear about
							\item[1.e] with \he{'El} to listen to someone
							\item[1.f] with \he{me'et yhwh} to hear something from Yahweh
							\item[1.g] with \he{p*iy} to hear that
							\item[1.h] with \he{le'mor} a direct question
							\item[2] to \textbf{listen to}
							\item[2.a] with accusative of the thing
							\item[2.b] with accusative of the person, to attend to someone's words carefully, or alternatively to hearken to someone
							\item[3] to \textbf{hear and accept a request}
							\item[4.a] to hear $\rightarrow$ to listen to, meaning to hearken $\rightarrow$ to \textbf{obey}...used absolutely, to be obedient \hyperref[1SAMUEL-15-22]{1 Samuel 15:22} \hyperref[2KINGS-14-11]{2 Kings 14:11} and elsewhere.
							\item[4.b] with \he{'El} of the person to satisfy, obey
								\begin{compactitem}
									\item A few other possible meanings are listed here.
								\end{compactitem}
							\item[5] to hear, meaning to \textbf{understand}
							\item[5.a] with accusative of the thing
							\item[5.b] followed by \he{mah}
							\item[5.c] used absolutely
							\item[5.d] \he{leb +some`a} a listening heart \hyperref[1KINGS-3-9]{1 Kings 3:9} (compare \hyperref[1KINGS-3-12]{verse 12}), a familiar expression in Egyptian wisdom literature, and probably taken over from there
							\item[6] special instances; within the sense of hearing is also included
							\item[6.a] to appreciate with one's spiritual hearing
							\item[6.b] in judicial proceedings
							\item[6.c] \he{+sAma` ha.tOb w:hArA`} \hyperref[2SAMUEL-14-17]{2 Samuel 14:17} to hear out, to sense what is good and what is bad...NRSV: discerning good and evil; REB: decide between right and wrong; a looser rendering of \he{+sAma`} here would be to decide
							\item[7.a] (isolated special instances)
						\end{compactitem}
					% \item[BDB , PDF ] \textbf{}
						% \begin{compactitem}
							% \item[1]
						% \end{compactitem}
					% \item[DCH PDF ] \textbf{} % don't include the actual page number, they repeat from volume to volume
						% \begin{compactitem}
							% \item[1]
						% \end{compactitem}
				\end{compactdesc}
				
				\begin{compactdesc}
					\item[Some OT uses] \textbf{}
						\begin{compactdesc}
							\item[{\hyperref[1SAMUEL-15-22]{1 Samuel 15:22}}]
						\end{compactdesc}
				\end{compactdesc}
				
		% --
		\paragraph{Dictionaries, Lexicons, and Encyclopedias}
		% --
				
			%
			\subparagraph{TWOT: \texorpdfstring{\he{+sAma`}}{} (938--939; \#2412)}
			%
				\label{SHAMA-TWOT}
				
				``The verb \he{+sAma`} is used 1050 times in the Qal, Niphal, Piel (twice), and Hiphil. Cognates are found in Akkadian, Aramaic, Arabic, Ugaritic, and Ethiopic. The basic idea is that of perceiving a message or merely a sound...
				
				``\he{+sAma`} has the basic meaning `to hear.' This is extended in various ways, generally involving an effective hearing or listening: 1) `listen to,' `pay attention,' 2) `obey' (with words such as `commandment' etc.), 3) `answer prayer,' `hear,' 4) `understand' and 5) `hear critically,' `examine (in court).' The derived stems have appropriately modified meanings.
				
				``Instances of the basic use of the verb are numerous. Examples are Num 12:2 (the Lord heard Miriam's and Aaron's grumbling), Deut 4:12 (the Israelites heard the sound of God's voice but saw no form) and Gen 3:8 (Adam and Eve heard the voice of God in the garden). The object of the hearing may be expressed in a dependent clause, such as in Gen 37:17 `heard (them) saying,' Gen 14: 14 `heard that his brother had been taken captive,' Jud 7:11 `hear what they say.'
				
				``The meaning `listen to' is illustrated in Gen 3:17 (Adam, to the voice of his wife, i.e. he followed her lead), ` Kgs 22:19, `Hear thou the word of the Lord,' Ps 81:11 [H 12]. `My people would not listen to my voice' and Prov 12: 15, `He that listens to counsel is wise.' This usage shades into that of `to obey,' such as in Ex 24:7; Isa 42:24 `obey his law': Neh 9:16, `They did not obey Thy commandments'; and Jer 35:18, `You have obeyed the commandment of Jonadab'...
				
				``In connection with answered prayer, God states a very important and sobering principle in Jer 11:10--11. Because Israel has refused to listen to the words of God when he spoke to them, they will find that when they cry to him in time of trouble he will not hear (respond to) their cry. Micah 7:7 expresses the confidence of the righteous, the one who himself heeds the voice of God, that God will indeed hear his prayer. A further strong word of encouragement is given us in Ps 94:9, `He who planted the ear, will he not hear?'
				
				``Effective hearing involves also the idea of `understanding.' Thus in Gen 11:7, after the confusion of languages at Babel, men could no longer `hear' (i.e. `understand') one another.''
				
			%
			\subparagraph{TDOT: \texorpdfstring{\he{+sAma`}}{} (15:253--279)}
			%
				\label{SHAMA-TDOT}
				
				``Social interaction takes place primarily in the form of linguistic communication; without language, speech, and hearing, the constitution of human society is inconceivable. This explains why suspension of hearing or refusal to hear is associated with chaos...for the OT we may cite Gen. 11:7: Yahweh dissolves the human society that is attempting to build the tower of Babel by making communication impossible: the people can no longer understand one another (\he{+sAma`}; the verb also has this meaning elsewhere, e.g., Gen 42:23; Dt. 28:49; 2 K. 18:26//Isa. 36:11; Jer. 5:15; Ezk. 3:6).
				
				``Dt 21:18--21 decrees the death penalty for someone who refused to listen to his parents (\he{+sAma`}, twice in v. 18); the pitiless harshness of the sentence points to the antisocial nature of the offense, which endangers society...
				
				``Hearing is critical for the interaction between God and human beings; the medium through which God makes his will known among his people (in commandments or mediated by the prophets) is the audible word...to exaggerate the point: one who cannot hear does not exist; one who can no longer hear, no longer communicate, is doomed (Ps. 38:14--15 [Eng. 13--14]).
			
		% --
		\paragraph{Septuagint}
		% --
			\label{OBEDIENCE-LXX}
			
			%
			\subparagraph{\texorpdfstring{\gr{ἀκοή}}{}}
			%
				\label{AKOE}
			
				\begin{compactdesc}
					\item[GELS 21a, PDF 62] \textbf{}
						\begin{compactitem}
							\item[1.a] \textit{that which is heard}: `message'...`report'...`rumour'
							\item[1.b] \textit{that which is publicly announced, proclaimed}
							\item[1.c] verbal noun of \gr{ἀκούω}
								\begin{compactitem}
									\item The use in 1 Kings (\hyperref[1SAMUEL-15-22]{1 Samuel 15:22}) is listed here
								\end{compactitem}
							\item[2] \textit{ear}: sense organ
						\end{compactitem}
					\item[GELSLEH 21] \textbf{}
						\begin{compactitem}
							\item \textit{sound}; \textit{report, tidings, news}; \textit{ear}; \textit{obedience}
								\begin{compactitem}
									\item The ``obedience'' meaning lists \hyperref[1SAMUEL-15-22]{1 Samuel 15:22}.
								\end{compactitem}
						\end{compactitem}
				\end{compactdesc}
				
				\begin{tabular}{ll}
					Total occurrences in the OT & 51
				\end{tabular}
				
				\begin{compactdesc}
					\item[Some OT uses] \textbf{}
						\begin{compactdesc}
							\item[{\hyperref[1SAMUEL-15-22]{1 Samuel 15:22}}]
						\end{compactdesc}
				\end{compactdesc}
		
	% -----
	\subsubsection{New Testament} % *NT
	% -----
		\label{OBEDIENCE-NT}
	
		% --
		\paragraph{KJV}
		% --
			\label{OBEDIENCE-NT-KJV}		
	
			% \begin{tabular}{ll}
				% Total occurrences in the NT & 
			% \end{tabular}
			
			\begin{compactdesc}
				\item[{\hyperref[2CORINTHIANS-10-6]{2 Corinthians 10:6}}]
			\end{compactdesc}
			
		% % --
		% \paragraph{Greek New Testament} % *GREEK
		% % --	
			% \label{OBEDIENCE-GREEK}
		
			% %
			% \subparagraph{\texorpdfstring{\gr{sample word}}{}}
			% %
				% \label{KATALLAGE}
			
				% \begin{compactdesc}
					% \item[BDAG , PDF ] \textbf{}
						% \begin{compactitem}
							% \item 
						% \end{compactitem}
					% \item[CGL , PDF ] \textbf{}
						% \begin{compactitem}
							% \item 
						% \end{compactitem}
					% \item[LSJ , PDF ] \textbf{}
						% \begin{compactitem}
							% \item 
						% \end{compactitem}
				% \end{compactdesc}
				
				% \begin{tabular}{ll}
					% Total occurrences in the NT & 
				% \end{tabular}
				
				% \begin{compactdesc}
					% \item[Some NT uses] \textbf{}
						% \begin{compactdesc}
							% \item[]
						% \end{compactdesc}
				% \end{compactdesc}
				
			% %
			% \subparagraph{TDNT: \texorpdfstring{\gr{sample word}}{}  (3:536--556)}
			% %
				% \label{KATALLAGE-TDNT}
		
	% % -----
	% \subsubsection{Book of Mormon} % *BOM
	% % -----
		% \label{OBEDIENCE-BOM}
	
		% \begin{tabular}{ll}
			% Total occurrences in the BOM & 
		% \end{tabular}
		
		% \begin{compactdesc}
			% \item[]
		% \end{compactdesc}
		
	% -----
	\subsubsection{Doctrine and Covenants} % *DC
	% -----
		\label{OBEDIENCE-DC}
	
		% \begin{tabular}{ll}
			% Total occurrences in the DC & 
		% \end{tabular}
		
		\begin{compactdesc}
			\item[{\hyperref[DC64-34]{DC 64:34}}]
		\end{compactdesc}
		
	% % -----
	% \subsubsection{Pearl of Great Price} % *PGP
	% % -----
		% \label{OBEDIENCE-PGP}
	
		% \begin{tabular}{ll}
			% Total occurrences in the PGP & 
		% \end{tabular}
		
		% \begin{compactdesc}
			% \item[]
		% \end{compactdesc}
		
% % ----------------------
% \subsection{Key Phrases} % *KEYPHRASES *PHRASES
% % ----------------------
	% \label{OBEDIENCE-PHRASES}

	% \begin{compactdesc}
		% \item[] \textbf{\#times} | list
			% % \begin{compactdesc}
				% % \item[Bible anchor verses] \phantom{.}
					% % \begin{compactdesc}
						% % \item[Romans 5:5] \gr{}
					% % \end{compactdesc}
				% % \item[Commentary] ---
			% % \end{compactdesc}
	% \end{compactdesc}
	
% % ----------------------
% \subsection{Bible Dictionaries} % *DICTIONARIES
% % ----------------------
	% \label{OBEDIENCE-BD}

% % ----------------------
% \subsection{Restoration Usage} % *RESTORATION
% % ----------------------
	% \label{OBEDIENCE-RESTORATION}

	% % -----
	% \subsubsection{Joseph Smith} % *JS
	% % -----
		% \label{OBEDIENCE-JS}
	
		% \begin{compactdesc}
			% \item[{\hyperref[KEY]{Date/Source}}]
		% \end{compactdesc}
		
	% % -----
	% \subsubsection{Temple Ordinances}
	% % -----
		% \label{OBEDIENCE-TEMPLE}
	
		% \begin{compactdesc}
			% \item[{\hyperref[KEY]{Date/Source}}]
		% \end{compactdesc}
	
	% % -----
	% \subsubsection{Brigham Young} % *BY
	% % -----
		% \label{OBEDIENCE-BY}
	
		% \begin{compactdesc}
			% \item[{\hyperref[KEY]{Date/Source}}]
		% \end{compactdesc}
	
% % ----------------------
% \subsection{Commentary and Studies} % *COMMENTARY *STUDIES
% % ----------------------
	% \label{OBEDIENCE-COMMENTARY}

	% % -----
	% \subsubsection{Key Points}
	% % -----
		% \label{OBEDIENCE-KEYS}
	
		% \begin{compactitem}
			% \item
		% \end{compactitem}
		
	% % -----
	% \subsubsection{Sample Study}
	% % -----
		% \label{OBEDIENCE-study}